\documentclass[11pt,a4paper]{article}

% ── Codificación y fuentes ────────────────────────────────────────────────────
\usepackage[utf8]{inputenc}
\usepackage[T1]{fontenc}
\usepackage{lmodern}
\usepackage[spanish]{babel}

% ── Geometría y espaciado ─────────────────────────────────────────────────────
\usepackage[top=2.5cm, bottom=2.5cm, left=2.8cm, right=2.8cm]{geometry}
\usepackage{setspace}
\setstretch{1.15}
\usepackage{parskip}

% ── Tablas ────────────────────────────────────────────────────────────────────
\usepackage{booktabs}
\usepackage{tabularx}
\usepackage{array}
\usepackage{longtable}
\usepackage{enumitem}

% ── Colores y cajas ───────────────────────────────────────────────────────────
\usepackage[dvipsnames]{xcolor}
\usepackage{tcolorbox}
\tcbuselibrary{skins, breakable}

% ── Cabeceras y pies de página ────────────────────────────────────────────────
\usepackage{fancyhdr}
\usepackage{lastpage}
\pagestyle{fancy}
\fancyhf{}
\fancyhead[L]{\small\textbf{Informe de Evaluación} --- Física para Bioanalistas}
\fancyhead[R]{\small Luisielis Marin}
\fancyfoot[C]{\small Página \thepage\ de \pageref{LastPage}}
\renewcommand{\headrulewidth}{0.4pt}

% ── Hiperenlaces ──────────────────────────────────────────────────────────────
\usepackage[colorlinks=true, linkcolor=NavyBlue, urlcolor=NavyBlue]{hyperref}

% ── Paleta de colores personalizados ─────────────────────────────────────────
\definecolor{desempeno_excelente}{HTML}{1A6B2F}
\definecolor{desempeno_bueno}{HTML}{2E5A9C}
\definecolor{desempeno_suficiente}{HTML}{8B6914}
\definecolor{desempeno_deficiente}{HTML}{C0392B}
\definecolor{desempeno_insuficiente}{HTML}{7B241C}
\definecolor{grisFondo}{HTML}{F5F5F5}
\definecolor{azulTitulo}{HTML}{1B2A6B}
\definecolor{verdeFortaleza}{HTML}{1A6B2F}
\definecolor{naranjaMejora}{HTML}{A04000}

% ── Estilos de cajas ─────────────────────────────────────────────────────────
\tcbset{
  cajaTitulo/.style={
    enhanced, breakable,
    colback=azulTitulo!8, colframe=azulTitulo,
    fonttitle=\bfseries\large, coltitle=white,
    attach boxed title to top left={yshift=-2mm, xshift=4mm},
    boxed title style={colback=azulTitulo},
    top=6pt, bottom=6pt, left=8pt, right=8pt,
  },
  cajaFortaleza/.style={
    enhanced, breakable,
    colback=verdeFortaleza!6, colframe=verdeFortaleza!60,
    fonttitle=\bfseries, coltitle=verdeFortaleza,
    top=4pt, bottom=4pt, left=8pt, right=8pt,
  },
  cajaMejora/.style={
    enhanced, breakable,
    colback=naranjaMejora!6, colframe=naranjaMejora!60,
    fonttitle=\bfseries, coltitle=naranjaMejora,
    top=4pt, bottom=4pt, left=8pt, right=8pt,
  },
  cajaRecomendacion/.style={
    enhanced, breakable,
    colback=grisFondo, colframe=gray!50,
    fonttitle=\bfseries, coltitle=black,
    top=4pt, bottom=4pt, left=8pt, right=8pt,
  },
}

% ─────────────────────────────────────────────────────────────────────────────
\begin{document}

% ══ PORTADA ══════════════════════════════════════════════════════════════════
\begin{center}
  {\Large\bfseries\color{azulTitulo} Universidad de Oriente/Núcleo Sucre --- Física para Bioanalistas}\\[4pt]
  {\large Informe de Evaluación Preliminar}\\[12pt]
  \rule{\linewidth}{1.2pt}\\[6pt]
  {\LARGE\bfseries Luisielis Marin}\\[4pt]
  \rule{\linewidth}{0.6pt}
\end{center}

% ══ FICHA TÉCNICA ════════════════════════════════════════════════════════════
\vspace{4pt}
\begin{tcolorbox}[cajaTitulo, title=Datos del informe]
\begin{tabular}{@{}llll@{}}
  \textbf{Estudiante:}  & Luisielis Marin  &
  \textbf{Fecha:}       & 2026-08-08 \\[4pt]
  \textbf{Archivo PDF:} & \texttt{Luisielis\_Marin.pdf}  &
  \textbf{Páginas:}     & 8 \\
\end{tabular}
\end{tcolorbox}

% ══ RESULTADO GLOBAL ═════════════════════════════════════════════════════════
\vspace{6pt}
\begin{center}
  \begin{tcolorbox}[
    enhanced,
    colback=white, colframe=desempeno_bueno,
    width=0.62\linewidth,
    boxrule=2pt,
    halign=center, valign=center,
    top=8pt, bottom=8pt,
  ]
    \centering
    {\large\bfseries Puntaje sugerido}\\[6pt]
    {\Huge\bfseries\color{desempeno_bueno} 7.6/10}\\[6pt]
    {\large Nivel de desempeño: \textbf{\color{desempeno_bueno} Bueno}}
  \end{tcolorbox}
\end{center}

% ══ RESUMEN GENERAL ══════════════════════════════════════════════════════════
\section*{Resumen general}
El estudiante demuestra una sólida comprensión de los principios físicos fundamentales de la unidad de líquidos y fenómenos de transporte, así como su aplicación en contextos bioanalíticos y clínicos. Las explicaciones son generalmente claras, detalladas y bien justificadas, especialmente en las preguntas 2, 3 y 4. Los cálculos son precisos y las unidades se manejan correctamente. La principal área de mejora es la omisión completa de la Pregunta 5 y una respuesta incompleta en una subsección de la Pregunta 1.

% ══ FORTALEZAS Y ÁREAS DE MEJORA ════════════════════════════════════════════
\vspace{4pt}
\begin{minipage}[t]{0.48\linewidth}
  \begin{tcolorbox}[cajaFortaleza, title={Fortalezas}]
\begin{itemize}[leftmargin=1.5em, itemsep=2pt]
  \item Profunda comprensión conceptual de la Ley de Laplace y el efecto del surfactante en la mecánica respiratoria.
  \item Excelente aplicación de la Primera Ley de Fick en el contexto de la hemodiálisis, incluyendo análisis de eficiencia y riesgos.
  \item Dominio de los principios de ósmosis y sus implicaciones clínicas (hemólisis, crenación) con explicaciones detalladas.
  \item Habilidad para realizar cálculos precisos y manejar unidades correctamente en todos los problemas numéricos resueltos.
  \item Capacidad para argumentar y justificar las respuestas con base en principios físicos, cumpliendo con el criterio fundamental de evaluación.
\end{itemize}
  \end{tcolorbox}
\end{minipage}
\hfill
\begin{minipage}[t]{0.48\linewidth}
  \begin{tcolorbox}[cajaMejora, title={Areas de mejora}]
\begin{itemize}[leftmargin=1.5em, itemsep=2pt]
  \item Completar todas las secciones del examen, ya que la Pregunta 5 fue omitida por completo.
  \item Profundizar en la comprensión de los factores que afectan la evaporación en condiciones extremas, como el 100\% de humedad relativa, y su impacto en la termorregulación.
\end{itemize}
  \end{tcolorbox}
\end{minipage}

% ══ OBSERVACIONES POR PREGUNTA ═══════════════════════════════════════════════
\section*{Observaciones por pregunta}

\begin{longtable}{>{\bfseries}p{2.8cm} p{1.6cm} p{7.0cm} p{2.8cm}}
  \toprule
  \textbf{Pregunta} & \textbf{Puntaje} & \textbf{Evaluación} & \textbf{Errores detectados} \\
  \midrule
  \endhead
  \bottomrule
  \endfoot
    \textbf{Pregunta 1: Termorregulación y Eficiencia de la Sudoración} & 1.6/2.0 & La parte (a) está excelentemente explicada, diferenciando claramente la eficiencia de la sudoración por evaporación versus goteo y su relación con el calor latente de vaporización. La parte (c) presenta un cálculo correcto y bien justificado, incluyendo la conversión a kcal. Sin embargo, la parte (b) es incompleta; aunque el estudiante menciona la dependencia de la evaporación con la diferencia de presión parcial de vapor, no explica qué sucede específicamente con el 100\% de humedad relativa y su impacto en la termorregulación de ambos pacientes. & \textit{Respuesta incompleta en la parte (b) sobre el efecto del 100\% de humedad relativa.} \\
    \midrule
    \textbf{Pregunta 2: Mecánica Respiratoria y Ley de Laplace} & 2.0/2.0 & La respuesta es sobresaliente en todas sus partes. La explicación de la Ley de Laplace y el colapso alveolar en (a) es precisa y concisa. La descripción del mecanismo de acción del surfactante en (b), incluyendo su dependencia del radio y cómo estabiliza las presiones alveolares, es muy completa y correcta. El cálculo en (c) es exacto y las unidades son adecuadas. & \textit{---} \\
    \midrule
    \textbf{Pregunta 3: Optimización en Hemodiálisis y Primera Ley de Fick} & 2.0/2.0 & La parte (a) demuestra una excelente aplicación de la Primera Ley de Fick, con una derivación algebraica clara y una conclusión correcta sobre el aumento del 40\% en la eficiencia. La parte (b) justifica de manera muy pertinente los riesgos mecánicos y estructurales de reducir excesivamente el espesor de la membrana. El cálculo en (c) es correcto y bien presentado, con unidades consistentes. & \textit{---} \\
    \midrule
    \textbf{Pregunta 4: Errores de Infusión Intravenosa y Ósmosis} & 2.0/2.0 & La respuesta es muy buena en todas sus secciones. La explicación de la hemólisis por agua destilada en (a) es detallada y conceptualmente sólida, cubriendo osmolaridad, flujo de solvente y fisiopatología. La descripción de la crenación por NaCl al 5\% en (b) es precisa y bien justificada. El cálculo de la presión osmótica en (c) es correcto y muestra un buen manejo de la fórmula y las unidades. & \textit{---} \\
    \midrule
    \textbf{Pregunta 5: Capilaridad y Toma de Muestras Sanguíneas} & 0.0/2.0 & La pregunta está en blanco. No hay respuesta del estudiante para ninguna de sus partes. & \textit{Pregunta sin responder.} \\
    \midrule
\end{longtable}

% ══ ERRORES DETECTADOS ═══════════════════════════════════════════════════════
\section*{Análisis de errores}

\subsection*{Errores conceptuales}
\begin{itemize}[leftmargin=1.5em, itemsep=2pt]
  \item Omisión de la explicación completa del efecto de 100\% de humedad relativa en la evaporación y la termorregulación del cuerpo (Pregunta 1b).
\end{itemize}

\subsection*{Errores procedimentales}
\begin{itemize}[leftmargin=1.5em, itemsep=2pt]
  \item Ninguno significativo en las preguntas respondidas.
\end{itemize}

% ══ RECOMENDACIONES AL ESTUDIANTE ════════════════════════════════════════════
\section*{Retroalimentación para el estudiante}
\begin{tcolorbox}[cajaRecomendacion, title={Dirigido a: Luisielis Marin}]
Tu desempeño en este examen es muy bueno, demostrando una sólida comprensión de la mayoría de los temas evaluados. Tus explicaciones son detalladas y tus cálculos precisos. Para mejorar aún más, es crucial que revises y completes todas las preguntas del examen, ya que la omisión de la Pregunta 5 afectó tu puntaje final. En la Pregunta 1b, asegúrate de entender completamente el impacto de la humedad relativa del 100\% en la evaporación y la termorregulación. Continúa practicando la aplicación de los principios físicos a escenarios clínicos, y presta atención a la completitud de tus respuestas.
\end{tcolorbox}

% ══ NOTA PARA EL PROFESOR ════════════════════════════════════════════════════
\section*{Nota para el profesor}
\begin{tcolorbox}[
  enhanced, breakable,
  colback=yellow!8, colframe=yellow!60!black,
  fonttitle=\bfseries, coltitle=black,
  title={Observaciones para revisión manual},
  top=4pt, bottom=4pt, left=8pt, right=8pt,
]
El estudiante ha demostrado un nivel de comprensión muy bueno en las preguntas respondidas (1-4). Las explicaciones son claras y los cálculos correctos, cumpliendo con el criterio de argumentación y justificación. La principal penalización se debe a la omisión completa de la Pregunta 5 y una respuesta incompleta en la Pregunta 1b, donde no se aborda el impacto de la humedad relativa del 100\% en la evaporación. El resto de las respuestas son ejemplares en su argumentación y aplicación de principios físicos.
\end{tcolorbox}

% ══ PIE DE INFORME ════════════════════════════════════════════════════════════
\vspace{12pt}
\begin{center}
  \footnotesize\color{gray}
  Informe generado automáticamente por el sistema de evaluación con IA (gemini-2.5-flash) y soporte LaTex. \\
  Este documento es una evaluación \textbf{preliminar} para ser revisado por el profesor antes de ser entregado al 
  estudiante. \\
  Generado el 2026-08-08 \textbullet\ BIOANALISIS Sistema Académico de Evaluación.
  \end{center}

\end{document}
